\documentclass[10pt,a4paper]{amsart}
\usepackage[margin=19mm]{geometry}
\usepackage{amsmath,amssymb,mathtools}
\usepackage[scr=rsfs,cal=euler]{mathalfa}
\usepackage{xfrac}
\usepackage[unicode,hidelinks,bookmarks=false]{hyperref}
\newcommand{\ie}{i.e.\ }
\newcommand{\cf}{cf.\ }
\newcommand{\resp}{respectively\ }
\newcommand{\Subsection}{\subsection}
\providecommand{\nobreakdash}{\nobreak\hbox{-}}
\setlength{\emergencystretch}{2em}
\multlinegap0pt
\usepackage{bm}
\usepackage{bbm}
\usepackage{dsfont}
\usepackage{xspace}
\usepackage{cancel}
\usepackage{mathtools}
\usepackage{amsthm}
\makeatletter
\@ifundefined{theorem}{\newtheorem{theorem}{Theorem}}{}
\@ifundefined{lemma}{\newtheorem{lemma}[theorem]{Lemma}}{}
\makeatother
\makeatletter
\expandafter\ifx\csname Corollary\endcsname\relax
\newenvironment{Corollary}{\goodbreak \begin{corollary}\rmfamily}{\end{corollary}}
\fi
\expandafter\ifx\csname Definition\endcsname\relax
\newenvironment{Definition}{\goodbreak \begin{definition}\rmfamily}{\end{definition}}
\fi
\expandafter\ifx\csname Example\endcsname\relax
\newenvironment{Example}{\goodbreak \begin{example}\rmfamily}{\bend\end{example}}
\fi
\expandafter\ifx\csname Lemma\endcsname\relax
\newenvironment{Lemma}{\goodbreak \begin{lemma}\slfamily}{\end{lemma}}
\fi
\newcommand{\N}{\ensuremath{\mathbb{N}}}
\expandafter\ifx\csname Proposition\endcsname\relax
\newenvironment{Proposition}{\goodbreak \begin{proposition}\rmfamily}{\bend\end{proposition}}
\fi
\expandafter\ifx\csname Remark\endcsname\relax
\newenvironment{Remark}{\goodbreak \begin{remark}\rmfamily}{\bend\end{remark}}
\fi
\newcommand{\T}{\ensuremath{\mathbb{T}}}
\expandafter\ifx\csname Theorem\endcsname\relax
\newenvironment{Theorem}{\goodbreak \begin{theorem}\slfamily}{\end{theorem}}
\fi
\renewcommand{\d}{\,\mathrm{d}}
\expandafter\ifx\csname eq\endcsname\relax
\newenvironment{eq}{
	\begin{equation}
	\begin{aligned}	
}{	\end{aligned}
	\end{equation}
}
\fi
\makeatother
\title[Positive quadrature and mobile sampling of multivariate trig]{Positive quadrature and mobile sampling of multivariate trigonometric polynomials}
\author{Stefan Kunis}
\date{Source: arXiv:2608.11915v1 (CC BY 4.0)}
\begin{document}
\maketitle

\noindent\emph{Source and licence.} Positive quadrature and mobile sampling of multivariate trigonometric polynomials. Version arXiv:2608.11915v1 (CC BY 4.0), distributed under the Creative Commons Attribution 4.0 International licence, as recorded in the arXiv licence metadata for this exact version and shown on the abstract page https://arxiv.org/abs/2608.11915v1. Full rights notice: licenses/arxiv-2608.11915-CC-BY-4.0.txt.\par\smallskip

\noindent\textit{Editorial scope.} Counted unit: the Lemma whose source label is \texttt{thm:gamma2} in the article (source0.tex lines 669--681, source label \texttt{thm:gamma2}), delivered together with the article's own complete original proof of it (source0.tex lines 682--720, one \texttt{proof} environment of 39 lines). No mathematics is written, repaired or re-derived: statement and proof are the article's own text line for line. Required context retained uncounted: eq:gamma (lines 606--608). Every definition, display and label of those lines is retained unchanged. Omitted from this package and not redistributed: 1--605, 609--668, 721--1006. Nothing inside a retained excerpt is omitted or altered: every retained body is delivered exactly as the article's lines are written, so no omission receipt of any kind accompanies this package. Citations: the retained excerpts carry no citation command at all, so no bibliography entry of the article is transcribed and no reference is redistributed. Reference adaptation: none was needed and none was made; every reference inside the delivered unit resolves inside it, so no unresolved reference or citation is produced. Printed slips kept literal: no printed word, symbol, hypothesis or number of the counted statement or of its proof was altered. Layout and class: the article's own class is \texttt{scrartcl} with packages including a4wide, amsfonts, amsmath, amsthm, mathtools, amssymb, array, multicol, wrapfig, dsfont, inputenc, fontenc, caption, subcaption; the wrapper is the lane's pinned \texttt{amsart} editorial wrapper at 10pt on A4, which restates the theorem-like environments the retained text needs and adds the article's own macro definitions that the retained text needs (\texttt{\textbackslash N}, \texttt{\textbackslash T}, \texttt{\textbackslash d}), with extra wrapper packages (bm, bbm, dsfont, xspace, cancel, mathtools). The delivered numbering is the unit's own. Assets: no retained span contains a figure environment, a graphics-inclusion command, a PDF-page inclusion, a table or a TikZ picture; no image file is copied into this package, the package declares no asset dependency and no image file is redistributed with it. Licence: version arXiv:2608.11915v1 of this article is distributed under the Creative Commons Attribution 4.0 International licence, as recorded in the arXiv licence metadata for this exact version and shown on the abstract page https://arxiv.org/abs/2608.11915v1. A full rights notice is delivered at licenses/arxiv-2608.11915-CC-BY-4.0.txt. Characters: every retained excerpt is pure ASCII, so the delivered engine, XeLaTeX with its default Latin Modern text font, reports no missing character.
\begin{align}\label{eq:gamma}
    y\mapsto \gamma(y)=\gamma_{d,n}=\left(y,(2n+1)y,\hdots,(2n+1)^{d-1} y\right) \mod 1,
\end{align}

\begin{lemma}\label{thm:gamma2}
 Let $d,n\in\N$, $d\ge 2$, then the curve \eqref{eq:gamma} has length
 \begin{align*}
    (2n+1)^{d-1}\le
    \ell(\gamma)
    =\left(\frac{(2n+1)^{2d}-1}{(2n+1)^2-1}\right)^{\frac12}
    \le \frac{3}{\sqrt{8}} (2n+1)^{d-1}
\end{align*}
and covering radius
\begin{align*}
    \frac{1}{4n+4} \le \delta = \frac{(2n+1)^{d-2}}{2(2n+1)^{d-1}+2} \le \frac{1}{4n+2}.
\end{align*}
\end{lemma}

\begin{proof}
We compute the length of this curve by
\begin{align*}
    \ell(\gamma)
    =\int_{\gamma}\d\sigma(x)
    =\int_0^1 \|\dot\gamma(y)\|_2\d y
    =\left(\sum_{r=0}^{d-1} (2n+1)^{2r}\right)^{\frac12}
    =\left(\frac{(2n+1)^{2d}-1}{(2n+1)^2-1}\right)^{\frac12}.
\end{align*}
The lower bound follows by considering the last summand, $r=d-1$, in the next to last term.
The upper bound is obtained via
\begin{align*}
 \frac{(2n+1)^{2d}-1}{(2n+1)^2-1}
 =(2n+1)^{2d-2}\left(\frac{(2n+1)^{2}-(2n+1)^{2-d}}{(2n+1)^2-1}\right)
 \le (2n+1)^{2d-2}\left(\frac{(2n+1)^{2}}{(2n+1)^2-1}\right)
\end{align*}
and the very last term is decreasing with $n$ and bounded by $\frac{9}{8}$ for $n\ge 1$.

Regarding the covering radius, we start by proving the upper bound. The points
\begin{align}\label{eq:lattice}
 \gamma\left(\frac{k}{(2n+1)^{d}}\right),\qquad k=0,\hdots,(2n+1)^{d}-1,
\end{align}
form a rank-1 lattice in $\T^{d}$. Placing a cube (square for $d=2$) with radius $\frac{1}{2(2n+1)}$ at every point gives a partition of $\T^{d}$. The covering radius of the curve is upper bounded by the covering radius of this subset of points.

The same upper bound can be obtained, by considering the hyperplane $\T^{d-1}\times\{0\}$: The points on the curve with last coordinate being zero, i.e.,
\begin{align*}
 \gamma\left(\frac{k}{(2n+1)^{d-1}}\right),\qquad k=0,\hdots,(2n+1)^{d-1}-1,
\end{align*}
form the rank-1 lattice \eqref{eq:lattice} with $d$ replaced with $d-1$ in $\T^{d-1}\times\{0\}$.
Hence, placing a $(d-1)$-cube (line segment for $d=2$, square for $d=3$) with radius $\frac{1}{2(2n+1)}$ at every point gives a partition of $\T^{d-1}$. Now moving these $(d-1)$-dimensional cubes along the curve also covers $\T^d$.

During this movement, the $d$-cube placed at the origin also covers points in the hyperplane $\T^{d-1}\times\{0\}$ with $\ell^\infty$-distance larger than $\frac{1}{2(2n+1)}$. We replace the $d$-cube of radius $\frac{1}{2(2n+1)}$ by one with smaller radius $\delta$.
When moving the center point of the cube along the curve up to the last coordinate being $\delta$, this cube still intersects with the 
hyperplane $\T^{d-1}\times\{0\}$ and the first coordinate (as well as all others) are increased by (at least) $\delta/(2n+1)^{d-1}$. Hence, the covering radius fulfils
\begin{align*}
 \frac{1}{2(2n+1)}=\delta+\frac{1}{(2n+1)^{d-1}}\delta.
\end{align*}
Solving for $\delta$ yields the result.
\end{proof}

\end{document}
